\documentclass{article}
\usepackage{titlesec}

\usepackage[citestyle=authortitle-terse,backref,backend=bibtex]{biblatex}
\addbibresource{homotopy-theory.bib}

\makeatletter
\renewcommand\thesection{}
\renewcommand\thesubsection{\@arabic\c@section.\@arabic\c@subsection}
\makeatother
\usepackage{tocloft}
\renewcommand{\thesubsection}{\arabic{subsection}}


\usepackage[left=4cm, right=4cm]{geometry}
\usepackage{palatino,eulervm,dutchcal,xcolor}
\usepackage{graphicx,subcaption,float}
\usepackage{enumitem,parskip,multicol,array}
\newcolumntype{L}{>{$}c<{$}}%table in mathmode
\usepackage{amsthm,amssymb,amsmath,mathtools,thmtools,tikz,tikz-cd}
\usetikzlibrary{%
	matrix,%
	calc,%
	arrows,%
	shapes,
	decorations.markings,backgrounds,calc,intersections}
\tikzcdset{scale cd/.style={every label/.append style={scale=#1},
		cells={nodes={scale=#1}}}}
\usepackage[bookmarks,bookmarksopen,bookmarksdepth=3]{hyperref}
\hypersetup{%colores
	colorlinks=true,
	urlcolor=blue,
	linkcolor=magenta,
	citecolor=blue,
	filecolor=blue,
	urlbordercolor=white,
	linkbordercolor=white,
	citebordercolor=white,
	filebordercolor=white}
\usepackage{cleveref}

\definecolor{blue-violet}{rgb}{0.54, 0.17, 0.89}
\definecolor{azure}{rgb}{0.0, 0.5, 1.0}
\definecolor{green(ncs)}{rgb}{0.0, 0.62, 0.42}
\definecolor{forestgreen}{rgb}{0.13, 0.55, 0.13}
\definecolor{limegreen}{rgb}{0.2, 0.8, 0.2}
\definecolor{palatinateblue}{rgb}{0.15, 0.23, 0.89}
\definecolor{trueblue}{rgb}{0.0, 0.45, 0.81}
\definecolor{goldenyellow}{rgb}{1.0, 0.87, 0.0}
\definecolor{fashionfuchsia}{rgb}{0.96, 0.0, 0.63}
\definecolor{brightcerulean}{rgb}{0.11, 0.67, 0.84}
\definecolor{jonquil}{rgb}{0.98, 0.85, 0.37}
\definecolor{lavendermagenta}{rgb}{0.93, 0.51, 0.93}
\definecolor{peru}{rgb}{0.8, 0.52, 0.25}
\definecolor{persimmon}{rgb}{0.93, 0.35, 0.0}
\definecolor{persianred}{rgb}{0.8, 0.2, 0.2}
\definecolor{persianblue}{rgb}{0.11, 0.22, 0.73}
\definecolor{persiangreen}{rgb}{0.0, 0.65, 0.58}
\definecolor{persianyellow}{rgb}{0.9, 0.89, 0.0}



\declaretheoremstyle[headfont=\color{trueblue}\normalfont\bfseries,]{colored1}
\declaretheoremstyle[headfont=\color{forestgreen}\normalfont\bfseries,]{colored2}
\declaretheoremstyle[headfont=\color{peru}\normalfont\bfseries,]{colored3}
\declaretheoremstyle[headfont=\color{persiangreen}\normalfont\bfseries,]{colored4}
\declaretheoremstyle[headfont=\color{brightcerulean}\normalfont\bfseries,]{colored5}
\declaretheoremstyle[headfont=\color{lavendermagenta}\normalfont\bfseries,]{colored6}
\declaretheoremstyle[headfont=\color{blue-violet}\normalfont\bfseries,]{colored7}
\declaretheoremstyle[headfont=\color{green(ncs)}\normalfont\bfseries,]{colored8}
\declaretheoremstyle[headfont=\color{peru}\normalfont\bfseries,]{colored9}
\declaretheoremstyle[headfont=\color{persiangreen}\normalfont\bfseries,]{colored10}

\declaretheorem[style=colored1,name=Theorem]{thm}
\declaretheorem[style=colored2,numberlike=thm,name=proposition]{prop}
\declaretheorem[style=colored3,numbered=no,name=Lemma]{lemma}
\declaretheorem[style=colored4,numbered=no,name=Corollary]{coro}
\declaretheorem[style=colored5,numbered=no,name=Example]{example}
\declaretheorem[style=colored5,numbered=no,name=Examples]{exemplos}
\declaretheorem[style=colored7,numbered=no,name=Exercise]{exercise}
\declaretheorem[style=colored6,numbered=no,name=Remark]{remark}
\declaretheorem[style=colored8,numbered=no,name=Claim]{claim}
\declaretheorem[style=colored9,numbered=no,name=Definition]{defn}
\declaretheorem[style=colored10,numbered=no,name=Question]{question}

\numberwithin{equation}{section}

\newcommand{\A}{\mathbb{A}}
\newcommand{\R}{\mathbb{R}}
\newcommand{\Z}{\mathbb{Z}}
\newcommand{\N}{\mathbb{N}}
\newcommand{\C}{\mathbb{C}}
\newcommand{\Q}{\mathbb{Q}}
\newcommand{\D}{\mathbb{D}}
\renewcommand{\P}{\mathbb{P}}

\newcommand{\Ac}{\mathcal{A}}
\newcommand{\Bc}{\mathcal{B}}
\newcommand{\Cc}{\mathcal{C}}
\newcommand{\Dc}{\mathcal{D}}
\newcommand{\Ec}{\mathcal{E}}
\newcommand{\Fc}{\mathcal{F}}
\newcommand{\Gc}{\mathcal{G}}
\newcommand{\Hc}{\mathcal{H}}
\newcommand{\Ic}{\mathcal{I}}
\newcommand{\Jc}{\mathcal{J}}
\newcommand{\Kc}{\mathcal{K}}
\newcommand{\Lc}{\mathcal{L}}
\newcommand{\Mc}{\mathcal{M}}
\newcommand{\Nc}{\mathcal{N}}
\newcommand{\Oc}{\mathcal{O}}
\newcommand{\Pc}{\mathcal{P}}
\newcommand{\Qc}{\mathcal{Q}}
\newcommand{\Rc}{\mathcal{R}}
\newcommand{\Sc}{\mathcal{S}}
\newcommand{\Tc}{\mathcal{T}}
\newcommand{\Uc}{\mathcal{U}}
\newcommand{\Vc}{\mathcal{V}}
\newcommand{\Wc}{\mathcal{W}}
\newcommand{\Xc}{\mathcal{X}}
\newcommand{\Yc}{\mathcal{Y}}
\newcommand{\Zc}{\mathcal{Z}}

\DeclareMathOperator{\img}{img}
\DeclareMathOperator{\Arg}{Arg}
\DeclareMathOperator{\id}{id}
\DeclareMathOperator{\pt}{pt}
\DeclareMathOperator{\Alt}{Alt}
\DeclareMathOperator{\sgn}{sgn}
\DeclareMathOperator{\hTop}{h-Top}
\DeclareMathOperator{\supp}{supp}
\DeclareMathOperator{\Int}{Int}
\DeclareMathOperator{\Ob}{Ob}
\DeclareMathOperator{\Mor}{Mor}
\DeclareMathOperator{\Top}{Top}
\DeclareMathOperator{\Set}{Set}
\DeclareMathOperator{\CGWH}{CGWH}
\DeclareMathOperator{\Hom}{Hom}
\DeclareMathOperator{\Map}{Map}
\DeclareMathOperator{\Tot}{Tot}
\DeclareMathOperator{\op}{op}
\DeclareMathOperator{\ev}{ev}
\DeclareMathOperator{\hofib}{hofib}
\DeclareMathOperator{\rel}{rel}
\DeclareMathOperator{\Nat}{Nat}
\DeclareMathOperator{\Arr}{Arr}
\DeclareMathOperator{\Funct}{Funct}
\DeclareMathOperator{\colim}{colim}
\DeclareMathOperator{\GL}{GL}
\DeclareMathOperator{\SL}{SL}
\DeclareMathOperator{\SO}{SO}
\DeclareMathOperator{\U}{U}
\DeclareMathOperator{\SU}{SU}
\DeclareMathOperator{\Sp}{Sp}
\DeclareMathOperator{\M}{M}
\DeclareMathOperator{\Aut}{Aut}
\DeclareMathOperator{\PGL}{PGL}
\DeclareMathOperator{\PSL}{PSL}

\begin{document}
	{\Large cohomology ring of $\C P^n$}
	\begin{exercise}
		Show that
		\[H^\bullet(\C P^n)\cong \Z[\alpha]/(\alpha^{n+1})\]
		where $\alpha$ has degree 2.
	\end{exercise}
\begin{proof}
		The CW structure of $\C P^n$ consists of one cell for every even dimension. This gives us the following chain complex:
	\[\begin{tikzcd}
		\Z\arrow[r]&0\arrow[r]&\Z\arrow[r]&\ldots\arrow[r]&\Z,\qquad\text{if }n\text{ is even}
	\end{tikzcd}\]
	\[\begin{tikzcd}
		\Z\arrow[r]&0\arrow[r]&\Z\arrow[r]&\ldots\arrow[r]&0,\qquad\text{if }n\text{ is odd}
	\end{tikzcd}\]
	which yields the cohomology
	\[H^i(\C P^n)=\begin{cases}
		\Z,\quad&i=0,2,4,\ldots 2n\\
		0,\quad&i\text{ otherwise}
	\end{cases}\]
	so that
	\begin{align*}
		H^\bullet(\C P^n)&=H^0(\C P^n)\oplus H^2(\C P^n)\oplus\ldots\oplus H^{2n}(\C P^{n})
	\end{align*}
	This means that the underlying group of the cohomology ring is the same as that of
	\[\Z[\alpha]/(\alpha^{n+1})\]
	where $\alpha$ has degree 2. To show that these groups are also isomorphic as algebras we can use Poincar\'e duality as follows.
	
	Consider the case $n=2$, where we may immediately multiply the generator of second cohomology group with itself:
	\begin{align*}
		H^2(\C P^2)\times H^2(\C P^2)&\to H^4(\C P^2)\\
		(\alpha,\alpha)\qquad\quad&\mapsto\alpha\smile\alpha =\alpha^2
	\end{align*}
	By Poincar\'e duality this map is nondegenerate symmetric bilinear form, so it must map generator to a generator. {\color{blue-violet}The fact that the product of the generator in degree 2 is the generator of degree 4} yields an homomorphism
	\iffalse\begin{align*}
		\varphi:H^\bullet(\C P^n)&\to\Z[\alpha]/(\alpha^{n+1})\\
		H^2(\C P^n)\ni\alpha&\mapsto\alpha
	\end{align*}
	satisfying
	\[\varphi(\alpha\smile\alpha)=\varphi()\]\fi
	\begin{align*}
		\varphi:\Z[\alpha]&\to H^\bullet(\C P^n)\\
		\alpha&\mapsto\alpha\in H^2(\C P^n)
	\end{align*}
	with kernel $(\alpha^{n+1})$ as desired.
	
	Now the case of $\C P^3$ is:
	\begin{align*}
		H^2(\C P^3)\times H^4(\C P^3)&\to H^6(\C P^3)\\
		(\alpha,\alpha^2)\qquad\quad&\mapsto\alpha\smile\alpha^2=\alpha^3
	\end{align*}
	which also  maps generator to generator, producing the desired algebra isomorphism. Notice we have used the group isomorphism $H^4(\C P^3)\approx H^4(\C P^2)$ when denoting the generator of $H^4(\C P^3)$ as $\alpha^2$. Such an isomorphism is induced by inclusion $\C P^{n-1}\hookrightarrow\C P^n$ via relative cohomology exact sequence.
	
	The case for dimension $n$ follows by induction.
\end{proof}
\end{document}